\scanpage{138}\mspage{70}

En effet, \deleted{une telle structure $n$-polygonale est définie en termes} \inserted{prenons d'abord des structures} $n$-polygônales orientées, soit $\overrightarrow{\mathrm{Pol}}(E)$, alors on a

(33)\qquad Soit $E$ de cardinal $n$, alors $(\mathbb Z/n\mathbb Z)^*$ opère canoniquement sur $\overrightarrow{\mathrm{Pol}}(E)$.

En effet, un élément de $\unsure{\overline{\mathrm{Pol}}|\overrightarrow{\mathrm{Pol}}}(E)$ s'interprète comme
\[
\begin{aligned}
&\text{opération de }\mathbb Z/n\mathbb Z=\Lambda_n\text{ sur }E\\
&\text{faisant de }E\text{ un }\Lambda_n\text{-torseur}\\
&\text{(i.e. une ``droite affine sur }\Lambda_n\text{'')}
\end{aligned}
\]
donc $\mathrm{Aut}(\Lambda_n)$ opère sur l'ens de ces structures, par composition. Si on identifie un $\omega\in\overrightarrow{\mathrm{Pol}}(E)$ comme une permutation circulaire, on aura
\[
\begin{aligned}
&\deleted{\unsure{\omega}}\ \underbrace{T_{\nu}\omega}=\omega^\nu\\
&\qquad\text{effet sur }\omega\in\overrightarrow{\mathrm{Pol}}(E)\subset\unsure{\mathfrak S|\mathcal S}_E\\
&\qquad\text{de }T_\nu\text{ associé à }\nu\in(\mathbb Z/n\mathbb Z)^*
\end{aligned}\tag*{(34)}
\]
On a donc
\[
\left\{\begin{aligned}
&T_{-\nu}\omega=\omega^{-\nu}=(\omega^\nu)^{-1}=(T_\nu\omega)^{-1}\\
&\qquad\text{structure polygonale orientée}\\
&\qquad\text{opposée à }\deleted{\unsure{\omega}}T_\nu\omega\\
&T_\nu\omega^{-1}=(T_\nu\omega)^{-1}\qquad\text{i.e. }T_\nu\text{ commute}\\
&\qquad\qquad\deleted{\unsure{\text{à}}}\enspace\omega\longmapsto\omega^{-1}\\
&\qquad\qquad\deleted{\text{\unsure{passage}}}\quad\text{passage d'une str.}\\
&\qquad\qquad\text{de polyg. orient. à \unsure{l'opposée|l'opposé}}
\end{aligned}\right.\tag*{(35)}
\]

% archive sheet number: 137
